\documentclass{article}
% For anonymized submission, change [preprint] to the default neurips_2025.
\usepackage{neurips_2025}
\makeatletter
\@anonymousfalse
\makeatother
\usepackage[utf8]{inputenc}
\usepackage[T1]{fontenc}
\usepackage{hyperref}
\usepackage{url}
\usepackage{booktabs}
\usepackage{amsfonts}
\usepackage{amsmath}
\usepackage{amssymb}
\usepackage{microtype}
\usepackage{xcolor}
\usepackage{graphicx}
% This prevents Overleaf from failing if a figure file is missing.
\newcommand{\safeincludegraphics}[2][]{%
  \IfFileExists{#2}{\includegraphics[#1]{#2}}{%
    \fbox{\parbox{0.9\linewidth}{\centering Missing figure file: \texttt{#2}}}%
  }%
}
\title{Medical Image Classification for Early Melanoma Detection}
\author{%
  Bryan Torres \\
  Virginia Tech \\
  Department of Computer Science \\
  \texttt{bryantorres@vt.edu}
  \And
  Roman Chenoweth \\
  Virginia Tech \\
  Department of Computer Science \\
  \texttt{romanc22@vt.edu}
  \And
  Cian Teague \\
  Virginia Tech \\
  Department of Computer Science \\
  \texttt{cianteague@vt.edu}
}
\begin{document}
\maketitle
\begin{abstract}
Melanoma is a serious form of skin cancer where early detection can substantially improve outcomes.
This project studies binary classification of dermoscopic lesion images into benign and malignant
categories using transfer-learned convolutional neural networks. Using an ISIC-derived Kaggle
dataset with 13,879 total images, we built a complete image-classification pipeline including
stratified train/validation splitting, preprocessing, data augmentation, two-phase fine-tuning, and
evaluation across five pretrained CNN architectures. The final benchmark shows strong performance
across all models, with ResNet-50 achieving the best overall result: AUC-ROC of 0.9800, accuracy
of 0.9185, F1 score of 0.9154, and sensitivity of 0.9850 at 80\% specificity on a balanced held-out
test set. Overall, the project demonstrates that a carefully implemented transfer-learning pipeline
can produce strong melanoma-screening performance while also showing why reliable splitting,
clinical metrics, and responsible interpretation are essential for medical-image machine learning.
\end{abstract}
\section{Introduction and Motivation}
Melanoma is a high-risk skin cancer, and early detection is one of the most important factors in
improving patient outcomes. In practice, visual inspection of skin lesions can be difficult because
benign and malignant lesions may share similar visual patterns. Dermoscopic images also vary in
illumination, zoom, lesion size, skin tone, and artifact content such as hair, rulers, shadows, and
dark borders. These issues make melanoma detection a meaningful machine learning problem,
especially when the goal is to support screening rather than replace clinical judgment.
The motivation for this project is to study whether modern deep neural networks can learn useful
visual representations for distinguishing malignant melanoma from benign lesions. We frame the
system as a decision-support prototype. In this setting, the most important concern is not simply
maximizing accuracy, but reducing missed malignant cases while still maintaining a reasonable
specificity. This means the project has asymmetric error costs: a false negative is usually more
concerning than a false positive because it may represent a missed cancer case.
From a technical perspective, this project is a supervised binary image-classification task. The main
work involved building a reliable training pipeline, comparing multiple pretrained CNN backbones,
and evaluating the models with clinically meaningful metrics. The final project focuses on whether
transfer learning, image augmentation, and threshold-based evaluation can produce strong and
interpretable results on a balanced held-out dermoscopy test set.
\section{Problem Definition}
Let an input dermoscopic image be represented as
\[
x \in \mathbb{R}^{224 \times 224 \times 3},
\]
where each image is resized to 224 by 224 pixels with three RGB channels. Each image has a binary
label
\[
y \in \{0,1\},
\]
where $y=0$ denotes benign and $y=1$ denotes malignant. The goal is to learn a model
\[
f_{\theta}(x) \in [0,1]
\]
that estimates the probability that the lesion is malignant. Given a decision threshold $\tau$, the
predicted label is
\[
\hat{y} = \mathbb{I}[f_{\theta}(x) \ge \tau].
\]
The learning problem is therefore to find parameters $\theta$ that separate benign and malignant
lesions well on unseen images. Because this is a medical-screening style task, success is measured
using more than raw accuracy. We evaluate AUC-ROC, precision, recall, F1 score, Brier score, and
sensitivity at fixed specificity values. Sensitivity at fixed specificity is especially important because
it shows how many malignant lesions the model can detect while holding the benign false-positive
rate to a controlled level.
\section{Related Work and Survey}
Deep learning has become a common approach for skin lesion analysis because convolutional neural
networks can learn hierarchical visual features directly from image data. Esteva et al.~\cite{esteva2017}
showed that deep CNNs trained with transfer learning can achieve dermatologist-level classification
performance on curated skin lesion datasets. This paper is important for our project because it
supports the core idea of using ImageNet-pretrained networks and fine-tuning them for medical
image classification.
Haenssle et al.~\cite{haenssle2018} compared a deep CNN against dermatologists for dermoscopic
melanoma recognition. A useful part of that study is its clinical framing: the model is evaluated at
specific sensitivity and specificity operating points rather than only through overall accuracy. This
influenced our decision to report sensitivity at 80\% and 90\% specificity.
The HAM10000 dataset introduced by Tschandl et al.~\cite{tschandl2018ham10000} and the ISIC
challenge benchmark summarized by Codella et al.~\cite{codella2019isic} helped standardize skin
lesion research. These datasets made it easier to compare models, but they also highlight common
medical-image issues such as class imbalance, limited metadata, duplicate-like samples, and possible
dataset bias across imaging conditions and patient groups.
For model design, ResNet~\cite{he2016resnet}, DenseNet~\cite{huang2017densenet},
EfficientNet~\cite{tan2019efficientnet}, and MobileNetV2~\cite{sandler2018mobilenetv2} offer
different trade-offs. ResNet uses residual connections to train deeper networks, DenseNet reuses
features through dense connections, EfficientNet scales network depth, width, and resolution in a
balanced way, and MobileNetV2 is designed for lighter deployment. Since these architectures have
different strengths, our project compares all of them under the same data split, preprocessing,
augmentation, and evaluation procedure.
Our regularization and imbalance-handling choices are also based on prior work. Focal
loss~\cite{lin2017focal} reduces the influence of easy examples and focuses learning on harder cases.
MixUp~\cite{zhang2018mixup} creates soft training examples by interpolating between images and
labels, while CutOut~\cite{devries2017cutout} masks random regions of an image to prevent the model
from depending too heavily on one visual cue. These methods are useful in our project because
medical images can contain artifacts or local patterns that may cause overfitting.
\section{Proposed Method}
\subsection{Intuition}
Our method is based on the idea that pretrained CNNs already contain useful low-level and mid-level
visual features, such as edges, textures, colors, and shapes. Instead of training a deep network from
scratch, we use transfer learning to adapt these features to dermoscopic lesion classification. This is
helpful because medical datasets are usually smaller than natural-image datasets, so training from
scratch would be more likely to overfit.
The project also uses a controlled comparison. Rather than changing several parts of the pipeline at
once, we keep the classification head, data split, augmentation procedure, and evaluation process
consistent across architectures. This makes the comparison more meaningful because differences in
performance are more likely to come from the backbone itself rather than from unrelated training
choices.
\subsection{Dataset and Preprocessing}
We used the ``Skin Cancer: Malignant vs. Benign'' dataset derived from the ISIC archive and hosted
on Kaggle. The dataset is provided as a directory-structured image collection, with RGB dermoscopic
images stored in JPEG format. Every image is resized to $224 \times 224$ pixels so that it can be used
with standard ImageNet-pretrained architectures.
The packaged dataset contains 13,879 total images. The provided training directory contains 11,879
images, with 6,289 benign images and 5,590 malignant images. The held-out test directory contains
2,000 images, evenly balanced with 1,000 benign and 1,000 malignant images. From the provided
training directory, we created a stratified 90/10 train/validation split. This resulted in 10,692
training images and 1,187 validation images. The training split contains 5,661 benign and 5,031
malignant examples, while the validation split contains 628 benign and 559 malignant examples.
\begin{figure}[t]
    \centering
    \safeincludegraphics[width=0.85\linewidth]{class_distribution.png}
    \caption{Class distribution for the packaged dataset. The training set is mildly imbalanced, while
    the held-out test set is balanced.}
    \label{fig:classdist}
\end{figure}
One important implementation issue involved validation splitting. Our first version used automatic
directory loading utilities, but because filenames were grouped alphabetically by class, the validation
slice could accidentally contain mostly one class. This caused unreliable validation metrics and made
early stopping untrustworthy. To fix this, we manually collected files per class and created a
stratified shuffle split. This correction was one of the most important parts of the project because
the final benchmark depends on a valid validation set.
\subsection{Model Architectures}
We benchmarked five pretrained CNN architectures: EfficientNet-B0, EfficientNet-B2, ResNet-50,
DenseNet-121, and MobileNetV2. Each model used the same classification head:
\[
\texttt{GlobalAveragePooling}
\rightarrow
\texttt{Dense(256, ReLU)}
\rightarrow
\texttt{BatchNorm}
\rightarrow
\texttt{Dropout(0.3)}
\rightarrow
\texttt{Dense(1, sigmoid)}.
\]
The sigmoid output gives the predicted probability that the lesion is malignant.
All architectures were trained using the same two-phase fine-tuning strategy. In the first phase, the
pretrained backbone was frozen and only the classification head was trained. In the second phase, the
backbone was unfrozen and fine-tuned end-to-end using a lower learning rate. Batch normalization
layers remained frozen during fine-tuning to preserve stable pretrained statistics.
\subsection{Losses, Augmentation, and Regularization}
The main benchmark used weighted binary cross-entropy. For a predicted probability $p=f_\theta(x)$
and label $y$, the loss is
\[
\mathcal{L}_{\mathrm{WCE}}
=
-
w_1 y \log p
-
w_0 (1-y)\log(1-p),
\]
where $w_0$ and $w_1$ are class weights derived from the training split. Since the dataset version is
only mildly imbalanced, the class weighting is not extreme, but it still helps account for the slightly
smaller malignant class.
We also considered focal loss as an imbalance-focused alternative:
\[
\mathcal{L}_{\mathrm{FL}}
=
-\alpha_t(1-p_t)^\gamma \log(p_t),
\]
where $p_t$ is the probability assigned to the true class, $\alpha_t$ is a class-balancing term, and
$\gamma$ controls how much the loss focuses on hard examples. In our final report, we treat focal
loss as an important extension and analysis point because increasing $\gamma$ can make the model pay
more attention to uncertain or misclassified images.
The augmentation pipeline included random horizontal and vertical flips, full-range random rotation,
random zoom, brightness and contrast perturbations, CutOut, and MixUp. MixUp creates soft training
examples:
\[
\tilde{x} = \lambda x_i + (1-\lambda)x_j,
\qquad
\tilde{y} = \lambda y_i + (1-\lambda)y_j.
\]
CutOut masks a random square region of the image. Together, these augmentations encourage the
model to learn more general lesion patterns instead of memorizing artifacts or small local details.
\begin{figure*}[t]
    \centering
    \safeincludegraphics[width=\textwidth]{augmentation_examples.png}
    \caption{Examples of the augmentation pipeline applied to a lesion image. MixUp is applied at
    the batch level, while CutOut masks a random region of an image during training.}
    \label{fig:augment}
\end{figure*}
\section{Experiments and Results}
\subsection{Experimental Questions and Testbed}
The experiments were designed to answer four main questions. First, can transfer-learned CNNs
achieve strong performance on benign versus malignant dermoscopic image classification? Second,
which pretrained architecture performs best under a consistent training and augmentation pipeline?
Third, do the models maintain high sensitivity at clinically meaningful specificity levels? Fourth,
what implementation issues affect the reliability of the benchmark?
All models were trained on the stratified training split and selected using validation loss. Final
evaluation was performed only on the balanced held-out test set. Predictions were exported as
probabilities and evaluated externally using scikit-learn. We computed accuracy, AUC-ROC,
precision, recall, F1 score, Brier score, and sensitivity at 80\% and 90\% specificity.
\subsection{Architecture Benchmark}
Table~\ref{tab:archbench} shows the final architecture benchmark. All five models achieved AUC
above 0.97, which suggests that transfer learning works well for this dataset. ResNet-50 had the best
overall performance, with the highest AUC, accuracy, F1 score, sensitivity at 90\% specificity, and
lowest Brier score.
\begin{table*}[t]
\centering
\small
\caption{Architecture benchmark on the balanced held-out test set. All models used the same
preprocessing, augmentation pipeline, two-phase fine-tuning procedure, and weighted cross-entropy
setup.}
\label{tab:archbench}
\begin{tabular}{lcccccccc}
\toprule
Architecture & Accuracy & AUC & Precision & Recall & F1 & Sens.@80\% Spec & Sens.@90\% Spec & Brier \\
\midrule
EfficientNet-B0 & 0.9090 & 0.9748 & 0.9495 & 0.8640 & 0.9047 & 0.9790 & 0.9480 & 0.0763 \\
EfficientNet-B2 & 0.9090 & 0.9743 & 0.9398 & 0.8740 & 0.9057 & 0.9830 & 0.9500 & 0.0768 \\
ResNet-50       & \textbf{0.9185} & \textbf{0.9800} & \textbf{0.9515} & 0.8820 & \textbf{0.9154} & \textbf{0.9850} & \textbf{0.9600} & \textbf{0.0666} \\
DenseNet-121    & 0.9130 & 0.9732 & 0.9366 & \textbf{0.8860} & 0.9106 & \textbf{0.9850} & 0.9450 & 0.0748 \\
MobileNetV2     & 0.9030 & 0.9713 & 0.9400 & 0.8610 & 0.8987 & 0.9830 & 0.9400 & 0.0808 \\
\bottomrule
\end{tabular}
\end{table*}
\begin{figure}[t]
    \centering
    \safeincludegraphics[width=\linewidth]{arch_benchmark.png}
    \caption{Comparison of AUC, recall, F1, and sensitivity at fixed specificity across the five
    benchmarked architectures.}
    \label{fig:archbench}
\end{figure}
\subsection{Observations and Interpretation}
The first main observation is that ResNet-50 was the strongest all-around model. It achieved an
AUC of 0.9800 and the lowest Brier score of 0.0666, which means it performed well both as a ranking
model and as a probability-producing model. Its sensitivity at 90\% specificity was also the best
among the models, which matters because melanoma screening needs high malignant-case detection
while still limiting unnecessary false positives.
The second observation is that DenseNet-121 was very competitive. It had the highest thresholded
recall at the default decision threshold, with recall of 0.8860 compared with ResNet-50's 0.8820.
This difference is small, but it shows that model selection depends on the chosen operating point.
If the main goal is default-threshold recall, DenseNet-121 is close to ResNet-50. If the main goal is
overall discrimination, calibration, and sensitivity at 90\% specificity, ResNet-50 is stronger.
The third observation is that MobileNetV2 performed slightly worse but still achieved an AUC of
0.9713. This makes it a reasonable deployment-friendly option if the project were extended into a
lightweight web demo or mobile setting. However, because this project focuses on classification
performance rather than deployment speed, ResNet-50 is the final preferred model.
The final observation is that the benchmark depends heavily on correct data handling. The validation
split issue showed that a model can appear unstable or misleading if the validation set is not properly
stratified. After correcting the split, the results became much more interpretable and consistent.
\subsection{Focal Loss Gamma Discussion}
The proposal feedback suggested that we analyze the focal loss focusing parameter $\gamma$. In this
project, $\gamma$ is important because it controls how strongly focal loss down-weights easy examples.
A smaller value behaves more like weighted binary cross-entropy, while a larger value focuses more
on hard or misclassified cases. In a melanoma-screening setting, this could be useful because difficult
malignant examples are clinically important.
For the final system, we report the weighted cross-entropy benchmark as the main result because it
was stable, easy to interpret, and already produced strong performance across all architectures.
Focal loss remains a useful extension for future work, especially if the dataset becomes more
imbalanced or if the goal shifts toward maximizing sensitivity at a stricter specificity level.
% Optional: If you have final focal-loss numbers, uncomment this table and fill it in.
%
% \begin{table}[t]
% \centering
% \small
% \caption{Optional focal-loss gamma comparison using the selected backbone. Fill this in only if you
% have final focal-loss results.}
% \label{tab:focalgamma}
% \begin{tabular}{lcccc}
% \toprule
% Loss Setting & AUC & Recall & F1 & Sens.@90\% Spec \\
% \midrule
% Focal Loss, $\gamma=0.5$ &  &  &  &  \\
% Focal Loss, $\gamma=1.0$ &  &  &  &  \\
% Focal Loss, $\gamma=2.0$ &  &  &  &  \\
% \bottomrule
% \end{tabular}
% \end{table}
\section{Software Package and Reproducibility}

The final project code, notebook workflow, and generated result files are maintained in a GitHub
repository so that the project can be reviewed, reproduced, and extended. The repository contains the
main exploration notebook, supporting Python code, saved benchmark outputs, and figures used in
the final report.

\subsection{Repository and Package Information}

\noindent\textbf{GitHub repository:} \\
\url{https://github.com/RomanLearns/melanoma} \\

% \noindent\textbf{Repository collaborator/access confirmation:} \\
% \texttt{[Fill in here once collaborator access is confirmed.]} \\

\noindent\textbf{Submitted package file:} \\
\texttt{[Fill in submitted tar.gz file name here.]} \\

\noindent\textbf{Notebook entry point:} \\
\texttt{capstone\_exploration.ipynb} \\

\noindent\textbf{Supporting code file:} \\
\texttt{cell\_fixes.py} \\

\subsection{Current Repository Structure}

The final project repository and submitted package are organized as follows:

\begin{verbatim}
melanoma/
  README.txt
  002.zip
  DOC/
    final_report.pdf
    final_report.tex
    graded_proposal.pdf
    graded_milestone.pdf
  SRC/
    capstone_exploration.ipynb
    cell_fixes.py
    capstone_results/
      arch_benchmark.pkl
      arch_benchmark.png
      augmentation_examples.png
      class_distribution.png
      sample_images.png
\end{verbatim}

The \texttt{README.txt} file serves as the user manual for the project and explains the project goal,
required packages, repository structure, dataset, dataset setup, and how to run the notebook workflow. The
\texttt{DOC/} directory contains the final report materials and copies of the graded proposal and
milestone documents. The \texttt{SRC/} directory contains the main project notebook,
\texttt{capstone\_exploration.ipynb}, the supporting implementation file \texttt{cell\_fixes.py}, and
the generated result files used in the report. The \texttt{capstone\_results/} directory stores the
saved benchmark results and figures, including the architecture benchmark plot, augmentation
examples, class distribution figure, and sample lesion images.
\section{Difficulties, Limitations, and Ethical Considerations}
The main technical difficulty was the validation split failure caused by directory-order loading.
This was an important lesson because the model itself was not the only possible source of error.
A flawed data pipeline can make validation results unreliable even if the architecture and training
code are correct. Fixing the split with a manual stratified shuffle was necessary before the benchmark
could be trusted.
A second difficulty was framework metric reliability. In repeated experiments, relying too heavily on
framework-side validation metrics can be risky because stateful metric objects may behave
unexpectedly across repeated model compilation. To reduce this risk, final metrics were computed
outside the training loop using scikit-learn.
This project also has medical-data limitations. The dataset is useful for an academic benchmark, but
we do not have enough metadata to evaluate fairness across skin tones, age groups, imaging devices,
or clinical sites. Because of this, the model should not be interpreted as a deployable clinical
diagnostic tool. It is a research prototype that demonstrates image-classification methods, not a
replacement for a dermatologist or medical professional.
Any future version of this system should include stronger external validation, fairness analysis,
clinical review, calibration testing, and clearer uncertainty reporting. A real medical deployment
would also require privacy safeguards, regulatory review, and careful human oversight.
\section{Conclusion and Future Work}
This project built a complete deep learning pipeline for binary melanoma image classification. We
prepared an ISIC-derived Kaggle dataset, corrected an important validation-splitting issue, applied
a consistent augmentation pipeline, trained five pretrained CNN architectures, and evaluated the
models on a balanced held-out test set. The strongest final result came from ResNet-50, which
achieved an AUC-ROC of 0.9800, accuracy of 0.9185, F1 score of 0.9154, and sensitivity of 0.9850
at 80\% specificity.
The main conclusion is that transfer learning is effective for this melanoma-classification benchmark,
but the reliability of the result depends heavily on careful preprocessing, stratified splitting, and
clinically meaningful evaluation metrics. The project also shows that different architectures offer
different trade-offs. ResNet-50 was the best overall model, DenseNet-121 was competitive for recall,
and MobileNetV2 remained useful as a lighter deployment candidate.
Future work could extend this project by running a more complete focal-loss and oversampling
ablation, adding Grad-CAM visualizations for interpretability, evaluating calibration more deeply,
and testing the model on external datasets. The most important future step would be external
validation across more diverse patient groups and imaging conditions before making any stronger
claims about real-world medical usefulness.
\begin{thebibliography}{99}
\bibitem{esteva2017}
A.~Esteva, B.~Kuprel, R.~A. Novoa, J.~Ko, S.~M. Swetter, H.~M. Blau, and
S.~Thrun.
\newblock Dermatologist-level classification of skin cancer with deep neural networks.
\newblock \emph{Nature}, 542:115--118, 2017.
\bibitem{haenssle2018}
H.~A. Haenssle, C.~Fink, R.~Schneiderbauer, \emph{et al.}
\newblock Man against machine: diagnostic performance of a deep learning convolutional neural
network for dermoscopic melanoma recognition in comparison to 58 dermatologists.
\newblock \emph{Annals of Oncology}, 29(8):1836--1842, 2018.
\bibitem{tschandl2018ham10000}
P.~Tschandl, C.~Rosendahl, and H.~Kittler.
\newblock The HAM10000 dataset: A large collection of multi-source dermatoscopic images of common
pigmented skin lesions.
\newblock \emph{Scientific Data}, 5:180161, 2018.
\bibitem{codella2019isic}
N.~Codella, V.~Rotemberg, P.~Tschandl, \emph{et al.}
\newblock Skin lesion analysis toward melanoma detection 2018: A challenge hosted by the
International Skin Imaging Collaboration (ISIC).
\newblock \emph{arXiv:1902.03368}, 2019.
\bibitem{he2016resnet}
K.~He, X.~Zhang, S.~Ren, and J.~Sun.
\newblock Deep residual learning for image recognition.
\newblock In \emph{CVPR}, 2016.
\bibitem{huang2017densenet}
G.~Huang, Z.~Liu, L.~van~der~Maaten, and K.~Q. Weinberger.
\newblock Densely connected convolutional networks.
\newblock In \emph{CVPR}, 2017.
\bibitem{tan2019efficientnet}
M.~Tan and Q.~V. Le.
\newblock EfficientNet: Rethinking model scaling for convolutional neural networks.
\newblock In \emph{ICML}, 2019.
\bibitem{sandler2018mobilenetv2}
M.~Sandler, A.~Howard, M.~Zhu, A.~Zhmoginov, and L.-C. Chen.
\newblock MobileNetV2: Inverted residuals and linear bottlenecks.
\newblock In \emph{CVPR}, 2018.
\bibitem{lin2017focal}
T.-Y. Lin, P.~Goyal, R.~Girshick, K.~He, and P.~Dollar.
\newblock Focal loss for dense object detection.
\newblock In \emph{ICCV}, 2017.
\bibitem{zhang2018mixup}
H.~Zhang, M.~Cisse, Y.~N. Dauphin, and D.~Lopez-Paz.
\newblock mixup: Beyond empirical risk minimization.
\newblock In \emph{ICLR}, 2018.
\bibitem{devries2017cutout}
T.~DeVries and G.~W. Taylor.
\newblock Improved regularization of convolutional neural networks with cutout.
\newblock \emph{arXiv:1708.04552}, 2017.
\end{thebibliography}
\end{document}